% !TEX root = 00_tmlr_beyond_scale.tex
%%%%%%%%%%%%%%%%%%%%%%%%%%%%%%%%%%%%%%%%%%%%%%%%%%%%%%%%%%%%%%%%%%%%%%%%%%%%%%%%
%  Beyond Scale — TMLR submission root
%
%  This root shares ALL section files (01–05, 99) and the bib with the DMLR
%  root 00_dmlr_beyond_scale.tex — single source of truth for content.
%  Only this file differs: tmlr.sty (auto-anonymizing in submission mode),
%  TMLR author block, mirrored abstract, Broader Impact as \section*.
%
%  SYNC NOTE: the abstract below is MIRRORED from 00_dmlr_beyond_scale.tex
%  (dmlr2e requires the abstract in its own root). If the abstract there is
%  edited — e.g. the model-count sentence under reconciliation — re-sync here.
%
%  Build: latexmk -pdf 00_tmlr_beyond_scale.tex
%  Anonymity: \usepackage{tmlr} (no option) hides authors + acknowledgments
%  are omitted below. For camera-ready: [accepted]; for arXiv: [preprint],
%  and re-enable \input{98_acknowledgments} (wrap its \acks in a plain
%  section — tmlr.sty does not define \acks).
%%%%%%%%%%%%%%%%%%%%%%%%%%%%%%%%%%%%%%%%%%%%%%%%%%%%%%%%%%%%%%%%%%%%%%%%%%%%%%%%

\documentclass[10pt]{article}
\usepackage{tmlr}
% If accepted: \usepackage[accepted]{tmlr}
% For preprint/arXiv: \usepackage[preprint]{tmlr}

\input{math_commands.tex}

\usepackage{hyperref}
\usepackage{url}
\usepackage{graphicx}
\usepackage{booktabs}
\usepackage{tabularx}
\usepackage{algorithm}
\usepackage{algorithmic}

\title{Beyond Scale: The Diversity Coefficient as a Data Quality Metric for Variability in Natural Language Data}

% Authors are hidden automatically while tmlr.sty is used without
% [accepted]/[preprint] — do not remove this block for submission.
\author{%
Brando Miranda\thanks{Correspondence to: \texttt{brando9@stanford.edu}\,\;and\;\texttt{sanmi@stanford.edu}}\space\space$^{1}$
\And
Alycia Lee$^{1}$
\And
Sudharsan Sundar$^{1}$
\AND
Allison Casasola$^{1}$
\And
Rylan Schaeffer$^{1}$
\And
Elyas Obadd$^{1}$
\AND
Sanmi Koyejo$^{1}$
\\[0.8em]
$^{1}$Department of Computer Science, Stanford University, Stanford, CA 94305, USA\\
\texttt{\{brando9, alylee15, sjsundar\}@stanford.edu}\\
\texttt{\{allyc, rschaef, eobbad, sanmi\}@stanford.edu}
}

\def\month{MM}  % camera-ready
\def\year{YYYY} % camera-ready
\def\openreview{\url{https://openreview.net/forum?id=XXXX}} % camera-ready

\begin{document}

\maketitle

\begin{abstract}
% SYNC: mirrored from 00_dmlr_beyond_scale.tex — keep identical.
Current trends in pre-training Large Language Models (LLMs) primarily focus on the scaling of model and dataset size.
While the \textit{quality} of pre-training data is considered an important factor for training powerful LLMs, it remains a nebulous concept that has not been rigorously characterized.
To this end, we propose a formalization of one key aspect of data quality -- measuring the \textit{variability} of natural language data -- specifically via a measure we call the diversity coefficient.
Our empirical analysis shows that the proposed diversity coefficient aligns with the intuitive properties of diversity and variability,
e.g., it increases as the number of latent concepts increases.
Then, we measure the diversity coefficient of publicly available pre-training datasets and demonstrate that their formal diversity is high compared to theoretical lower and upper bounds.
Finally, we conduct a comprehensive set of controlled \textit{interventional} experiments with GPT-2 and LLaMAv2 that demonstrate the diversity coefficient of pre-training data characterizes useful aspects of downstream model evaluation performance---totaling 44 models of various sizes (51M to 7B parameters).
We conclude that our formal notion of diversity is an important aspect of data quality that captures variability and causally leads to improved evaluation performance.
\end{abstract}

% algorithmic.sty re-defines \AND inside every algorithmic environment via
% \newcommand, which clashes with tmlr.sty's author-block \AND. The author box
% is fully processed at \maketitle above, so drop tmlr's \AND here; our
% pseudocode never uses the \AND keyword.
\let\AND\undefined

\input{01_introduction}

\input{02_method}

\input{03_experiments}

\input{04_related_work}

\input{05_discussion}

\section*{Broader Impact Statement}
This work proposes the diversity coefficient as a quantitative metric for data quality in LLM pre-training.
By making data diversity formally measurable, we enable researchers and practitioners to design more principled and transparent data curation pipelines.
A potential concern is that diversity metrics could be gamed or misused to justify data collection practices that raise privacy or copyright issues; we encourage responsible use alongside appropriate data governance.
All datasets used in this study are publicly available.

% Acknowledgments intentionally omitted in the anonymized submission build.
% On acceptance, re-enable via a plain \section*{Acknowledgments} (tmlr.sty
% has no \acks macro): see 98_acknowledgments.tex for the content.

\bibliographystyle{tmlr}
\bibliography{dmlr_bib_refs}

\appendix

\input{99_appendix}

\end{document}
